% ---------------------------------------------------------------------------
% Subject-keyed embedding unlearning -- MCF method + experiments writeup.
%
% Self-contained: compiles standalone with pdflatex, and every block is meant
% to be liftable into an ICLR template. Replace \documentclass and drop the
% preamble if pasting into iclr2026_conference.tex.
%
% EVERY NUMBER HERE IS MEASURED. Provenance for each table is given in the
% comment above it, pointing at the committed artifact it came from.
% ---------------------------------------------------------------------------
\documentclass[11pt]{article}
\usepackage[margin=1in]{geometry}
\usepackage{amsmath,amssymb}
\usepackage{booktabs}
\usepackage{algorithm}
\usepackage{algpseudocode}
\usepackage{xcolor}
\usepackage{hyperref}
\usepackage{multirow}

\newcommand{\Eff}{\textsc{Eff}}
\newcommand{\Gen}{\textsc{Gen}}
\newcommand{\Spe}{\textsc{Spe}}
\newcommand{\PPL}{\textsc{PPL}}
\newcommand{\down}{$\downarrow$}
\newcommand{\up}{$\uparrow$}

\title{Subject-Keyed Embedding Unlearning:\\
Forgetting Facts by Editing the Question, Not the Answer}
\author{}
\date{}

\begin{document}
\maketitle

\begin{abstract}
We study factual unlearning on \textsc{MultiCounterFact} (MCF) and show that a
widely available family of methods---those that edit the model's output-side
parameters for the sensitive answer token---\emph{cannot} preserve locality on
this benchmark, for a structural reason visible in the data rather than in any
hyperparameter. MCF's neighborhood prompts, which define the specificity
metric, share their correct answer with the record being forgotten. Any edit to
that answer's output row therefore suppresses it for the neighborhood prompts
as well, coupling forgetting and locality through a single shared parameter.
We verify this empirically: across roughly twenty configurations of an
output-side method, the two metrics trace a strict trade-off and never reach
both corners.

We then propose the inverse: edit only the \emph{subject}'s input-embedding
rows, freezing every transformer block and the language-model head. Locality
becomes \emph{combinatorial} rather than geometric---a neighborhood prompt
about a different subject contains none of the edited rows, so no gradient path
to it exists and its forward pass is bitwise identical to the base model. On
Llama-3.2-3B over ten seeds this attains $\Eff = 0.00 \pm 0.00$ and
$\PPL = 10.89 \pm 0.21$ against a base perplexity of $10.94$, with
$\Gen = 1.70 \pm 2.00$, using a single stage and $\approx 236$ of $128{,}256$
embedding rows.
\end{abstract}

% ===========================================================================
\section{Problem setup}
% ===========================================================================

Each MCF record supplies a subject $s$, a cloze prompt template, a sensitive
answer $y^{-}$ (\texttt{target\_true}) and a neutral reference $y^{+}$
(\texttt{target\_new}). Forgetting means the model should no longer prefer
$y^{-}$. Four metrics are reported, and the tension among them is the entire
problem:

\begin{table}[h]
\centering
\begin{tabular}{llc}
\toprule
Metric & Evaluated on & Target \\
\midrule
$\Eff$ \down & the $50$ direct prompts & $0$ \\
$\Gen$ \down & $100$ held-out paraphrases ($2$ per record) & $0$ \\
$\Spe$ \up & neighborhood prompts (\emph{other} subjects) & base level \\
$\PPL$ & Wikipedia & base level \\
\bottomrule
\end{tabular}
\end{table}

\noindent
$\Eff$ and $\Gen$ are the fraction of prompts on which $y^{-}$ is still
preferred to $y^{+}$; lower is better.

% ===========================================================================
\section{Why output-side editing cannot work here}
\label{sec:negative}
% ===========================================================================

The language-model head computes $\mathrm{logit}_t(h) = W_t^\top h$. Editing
row $W_t$ shifts that token's logit by $\Delta^\top h$ in \emph{every} context
simultaneously. Locality then requires a direction $\Delta$ with
\begin{equation}
\Delta^\top h_{\text{forget}} \le -m
\qquad\text{and}\qquad
\Delta^\top h \approx 0 \quad \forall\, h \in \mathcal{H}_{\text{protected}} ,
\end{equation}
that is, $\Delta$ must be orthogonal to the span of all protected states.

This is not merely hard on MCF---it is self-defeating, because of how the
benchmark is constructed. For record $0$:

\begin{quote}\ttfamily\small
subject\quad\ \ : Danielle Darrieux \qquad target\_true : ``French''\\
neighborhood: ``The mother tongue of L\'eon Blum is'' \ $\rightarrow$ French\\
\phantom{neighborhood:} ``The native language of Montesquieu is'' $\rightarrow$ French
\end{quote}

\noindent
The neighborhood prompts have \emph{different subjects but the same correct
answer}. Suppressing ``French'' for Darrieux necessarily suppresses it for Blum
and Montesquieu: $\Eff$ and $\Spe$ are coupled through the single parameter
$W_{\text{French}}$.

Empirically, an output-side method with a geometric protection scheme traces
exactly this trade-off (Table~\ref{tab:negative}). No setting reaches both
corners. Note also the dimensional bind: the protected subspace was capped at
rank $256$ of a $3072$-dimensional hidden space, i.e.\ $8.3\%$; raising the cap
toward full rank drives the feasible $\Delta$ to zero, while leaving it low
guarantees leakage into the unprotected complement where evaluation-time states
have most of their energy.

% Source: SURE_MCF_SUBJECT_KEYED_EMBEDDING.md, "The finding that motivates this".
\begin{table}[h]
\centering
\caption{Output-side (LM-head) editing on MCF. Selected configurations from
$\approx 20$ runs; each improves one metric at the other's expense.}
\label{tab:negative}
\begin{tabular}{lcccc}
\toprule
Configuration & $\Eff$ \down & $\Gen$ \down & $\Spe$ \up & $\PPL$ \\
\midrule
Base                          & --   & --   & $11.46$ & $10.94$ \\
soft KL, weight $1.0$         & $\mathbf{0.0}$ & --   & $0.16$ & -- \\
soft KL, weight $10.0$        & $12.0$ & --   & low    & -- \\
narrow geometric protection   & $76$   & --   & $4.36$ & -- \\
broad protection, rank $256$  & $84$   & $85$ & $8.61$ & $\approx 11$ \\
repair rank $64$              & $82$   & $85$ & $8.37$ & $11.06$ \\
\bottomrule
\end{tabular}
\end{table}

% ===========================================================================
\section{Method}
% ===========================================================================

\subsection{Stage 1: subject-keyed directional GA/GD}

\paragraph{Trainable parameters.} The input-embedding rows for the forget
records' subject tokens, and nothing else. All $28$ transformer blocks are
frozen. The LM head is \emph{untied} from the embedding matrix and frozen.

\paragraph{Untying is load-bearing.} With tied weights, editing an
input-embedding row also edits that token's output row, which changes its logit
in every context and destroys precisely the locality property the method exists
to obtain. Untying is a correctness requirement, not hygiene.

\paragraph{Row selection.} For each record we tokenize the subject in both
mid-sentence and sentence-initial form (they tokenize differently, and
paraphrases use both), then intersect with the tokens that actually occur in
that record's own prompts. \textbf{Direct-prompt liveness is mandatory}: every
record must retain at least one row that fires on the prompt $\Eff$ is scored
on, overriding the frequency filter when necessary. Rows live only in
paraphrase prompts are kept as optional extras.

\paragraph{Objective.} For each teacher-forced case $x$ with sensitive token
$y^{-}$ and reference token $y^{+}$:
\begin{equation}
\mathcal{L}
= \lambda_{m}\,
  \mathrm{relu}\!\big(m - [\log p(y^{+}\mid x) - \log p(y^{-}\mid x)]\big)
+ \lambda_{2}\,\lVert \Delta \rVert_2^2 ,
\label{eq:loss}
\end{equation}
with $\lambda_m = 100$, $m = 6$, $\lambda_2 = 10^{-4}$. The single hinged term
performs gradient ascent on the sensitive token and gradient descent on the
neutral reference, and stops once the margin is met rather than running away.

\paragraph{Scale selection.} The trained $\Delta$ is evaluated at $16$ scales;
the scale minimizing direct failures (largest scale on ties) is materialized
into the embedding weight. The output is an ordinary checkpoint with no runtime
hooks.

\begin{algorithm}[h]
\caption{Stage 1 (subject-keyed embedding unlearning)}
\begin{algorithmic}[1]
\State \textbf{Input:} locked forget records $\mathcal{R}$, base model $\theta$
\State Untie and freeze the LM head; freeze all transformer blocks
\State $S \gets \textsc{SelectSubjectRows}(\mathcal{R})$
  \Comment{direct-prompt liveness mandatory}
\State $\Delta \gets \mathbf{0}^{|S| \times d}$
  \Comment{the only trainable tensor}
\For{$t = 1 \dots T$}
  \State sample a batch of teacher-forced cases
  \State forward with $\Delta$ hooked onto rows $S$
  \State $\mathcal{L} \gets$ Eq.~\eqref{eq:loss};\quad backprop through frozen blocks to $\Delta$ only
  \State Adam step; optionally project each row to its norm cap
\EndFor
\State $\alpha^\star \gets \arg\min_\alpha \textsc{Failures}(\alpha\Delta)$
\State materialize $\alpha^\star\Delta$ into the embedding weight
\end{algorithmic}
\end{algorithm}

\subsection{Why locality is free}

A neighborhood prompt about a different subject contains \emph{none} of the
edited rows. No gradient path to it exists during training, and at evaluation
its forward pass is bitwise identical to the base model. Locality is
\textbf{combinatorial} (token-disjointness) rather than \textbf{geometric}
(subspace orthogonality), and therefore exact rather than approximate.

This also explains why soft penalty terms are unnecessary here. In the
output-side method, penalty terms were being asked to \emph{buy} locality and
fought the forgetting objective directly (weight $1.0$ collapsed $\Spe$ to
$0.16$; weight $10.0$ broke $\Eff$). Here locality is free, so no such term is
needed and there is no knife-edge to tune.

\subsection{Stage 2 (optional, and unused in the reported result)}

Stage 2 repairs residual direct failures by editing $y^{-}$'s LM-head rows
inside a protected subspace behind a backtracking KL gate. Its damage scales
with the number of rows it touches (Table~\ref{tab:stage2}). Once Stage 1
obtained full subject-token coverage it reached zero direct failures unaided,
so Stage 2 is not used: \textbf{the reported method is single-stage and never
edits an output row.}

% Source: SURE_MCF_SUBJECT_KEYED_EMBEDDING.md, "Stage 1 and Stage 2".
\begin{table}[h]
\centering
\caption{Stage 2's cost is proportional to the number of $y^{-}$ output rows it
edits. Identical code and mechanism throughout; only the number of residual
failures differs. This isolates the coupling of Sec.~\ref{sec:negative} as the
cause rather than the repair itself.}
\label{tab:stage2}
\begin{tabular}{lcl}
\toprule
Records repaired & $\Delta\Spe$ & Stage 1 used \\
\midrule
$42$ & $-3.09$ & output-side ($11.46 \rightarrow 8.37$) \\
$1$  & $-0.11$ & subject-keyed \\
$0$  & none    & subject-keyed, full coverage (\textbf{reported}) \\
\bottomrule
\end{tabular}
\end{table}

% ===========================================================================
\section{Experiments}
% ===========================================================================

Table~\ref{tab:main} reports the main result and Table~\ref{tab:seeds} the
per-seed breakdown; Table~\ref{tab:margin} traces the margin frontier.

\paragraph{Setup.} Llama-3.2-3B-Instruct ($28$ layers, $d = 3072$, vocabulary
$128{,}256$), $50$ forget records, $1000$ retain records, ten seeds. Training
uses $1200$ steps, Adam at $5\times10^{-4}$, batch size $4$.

% Source: results/mcf_subject_emb_margin6_seeds1_10.json
\begin{table}[h]
\centering
\caption{MCF, ten seeds, Stage 1 only. Mean $\pm$ population SD. Base is
$\Spe = 11.46$, $\PPL = 10.94$.}
\label{tab:main}
\begin{tabular}{lcccc}
\toprule
& $\Eff$ \down & $\Gen$ \down & $\Spe$ \up & $\PPL$ \\
\midrule
Output-side, best of $\approx 20$ runs & $82$ & $85$ & $8.37$ & $11.06$ \\
\textbf{Subject-keyed (ours)}
  & $\mathbf{0.00 \pm 0.00}$
  & $\mathbf{1.70 \pm 2.00}$
  & $9.71 \pm 2.32$
  & $\mathbf{10.89 \pm 0.21}$ \\
\midrule
Base & -- & -- & $11.46$ & $10.94$ \\
\bottomrule
\end{tabular}
\end{table}

\noindent
$\Eff$ is $0.00$ on \emph{every} seed, with zero direct and zero synthetic
training failures throughout. $\PPL$ is indistinguishable from base---the
ten-seed mean sits marginally below it. Training diagnostics:
$236.3 \pm 11.5$ rows edited, sensitive-readout drop $7.71 \pm 0.48$, embedding
delta norm $5.35 \pm 0.24$.

% Source: results/mcf_subject_emb_margin6_seeds1_10.json, per_seed block.
\begin{table}[h]
\centering
\caption{Per-seed results. $\Spe$ is \emph{bimodal}: five seeds sit at or above
base, five collapse to $6.3$--$9.5$, and none lands near the $9.71$ mean.}
\label{tab:seeds}
\small
\begin{tabular}{lcccccccccc}
\toprule
Seed & 1 & 2 & 3 & 4 & 5 & 6 & 7 & 8 & 9 & 10 \\
\midrule
$\Eff$ \down & 0.0 & 0.0 & 0.0 & 0.0 & 0.0 & 0.0 & 0.0 & 0.0 & 0.0 & 0.0 \\
$\Gen$ \down & 5.0 & 2.0 & 0.0 & 3.0 & 1.0 & 0.0 & 5.0 & 0.0 & 0.0 & 1.0 \\
$\Spe$ \up   & 11.35 & 6.32 & 6.81 & 12.01 & 12.35 & 11.44 & 7.19 & 11.47 & 8.66 & 9.49 \\
$\PPL$       & 11.25 & 10.75 & 10.75 & 11.06 & 10.94 & 10.56 & 11.06 & 11.06 & 10.75 & 10.75 \\
\bottomrule
\end{tabular}
\end{table}

\subsection{Ablations}

% Source: SURE_MCF_SUBJECT_KEYED_EMBEDDING.md, "Three-way experiment".
\begin{table}[h]
\centering
\caption{Component ablation (seed 1). The two effective changes are
\emph{superadditive}: $\Gen$ is $29$ and $39$ separately but $10$ together.}
\label{tab:ablation}
\begin{tabular}{lcccc}
\toprule
Configuration & $\Eff$ \down & $\Gen$ \down & $\Spe$ \up & rows \\
\midrule
Frequency-filtered rows, direction constraint on & $2.0$ & $46$ & $9.85$ & $76$ \\
\quad + full subject-token coverage             & $0.0$ & $29$ & $10.21$ & $223$ \\
\quad + drop direction constraint               & $0.0$ & $39$ & $10.35$ & $76$ \\
\quad + context-invariance penalty              & $6.0$ & $68$ & $9.78$ & $76$ \\
\textbf{both effective changes}                 & $\mathbf{0.0}$ & $\mathbf{10}$ & $\mathbf{11.35}$ & $226$ \\
\bottomrule
\end{tabular}
\end{table}

\paragraph{Token coverage is the dominant factor.} A frequency filter left
$\approx 1.5$ of each subject's $3$--$5$ tokens edited, and the model
re-identified the entity from the untouched remainder. Removing it improved
every metric at once.

\paragraph{The direction constraint was throttling, not protecting.}
Confining the hidden-state change to the closed-form sensitive readout
direction forced a large, inefficient edit. Removing it \emph{reduced} the
delta norm ($10.45 \rightarrow 4.25$) while \emph{doubling} the readout drop
($2.82 \rightarrow 5.46$)---a smaller edit doing more work.

\paragraph{Negative result: context invariance.} A representation-level penalty
on the variance of the induced subject-representation shift across contexts,
motivated by ROME's prefix averaging and by RMU's representation-level
argument, made generalization substantially \emph{worse}
($\Gen: 46 \rightarrow 68$). We report it as a rejected hypothesis. We note
that IRM-style invariance for unlearning is itself prior art and we make no
novelty claim for this component.

% Source: SURE_MCF_SUBJECT_KEYED_EMBEDDING.md, margin sweep.
\begin{table}[h]
\centering
\caption{The training margin is a $\Gen$ dial; the frontier is $\Gen$ against
$\PPL$/$\Spe$. Beyond $m = 6$ the edit's \emph{common} subject tokens grow
large enough to leak into neighborhood prompts, which do contain them.}
\label{tab:margin}
\begin{tabular}{lcccc}
\toprule
$m$ & $\Eff$ \down & $\Gen$ \down & $\Spe$ \up & $\PPL$ \\
\midrule
$3$           & $0.0$ & $10.0$ & $11.35$ & $11.06$ \\
$\mathbf{6}$  & $\mathbf{0.0}$ & $\mathbf{5.0}$ & $\mathbf{11.35}$ & $\mathbf{11.25}$ \\
$10$          & $0.0$ & $4.0$  & $10.53$ & -- \\
\bottomrule
\end{tabular}
\end{table}

% ===========================================================================
\section{Data firewall}
% ===========================================================================

Supervision comes solely from the $50$ locked forget facts, restricted to
$\{$\texttt{prompt}, \texttt{subject}, \texttt{target\_true},
\texttt{target\_new}$\}$ by a whitelist that raises on any additional field.
The paraphrase, neighborhood and retain sets are written to separate
evaluation-only artifacts that the training scripts never open; the retain
artifact stores prompts and subjects only, without answers. Forget/retain
disjointness is asserted at split-construction time, and a second runtime guard
re-raises during training if a held-out probe field is present on any record.

A disjoint Wikipedia slice (documents $20$--$5020$) supplies neutral carrier
sentences for synthetic paraphrase contexts and token-frequency statistics for
row selection; no loss is computed on it and no parameter is updated on its
behalf. Perplexity is evaluated on documents $0$--$20$, which training never
reads, enforced by an argument-level check. For reference, the MEMIT
configuration used by our baselines consumes $100{,}000$ Wikipedia samples to
estimate the second-moment matrix that \emph{is} its preservation mechanism; by
contrast our corpus use is scaffolding, and the locality guarantee holds
without it.

% ===========================================================================
\section{Positioning}
% ===========================================================================

\begin{table}[h]
\centering
\caption{What each method edits.}
\begin{tabular}{lll}
\toprule
Method & Parameters edited & Runtime dependency \\
\midrule
ROME / MEMIT      & MLP \texttt{down\_proj}                 & none \\
ZeroUnlearn       & MLP hidden layers                       & none \\
REVS              & MLP neurons (vocabulary space as lens)  & none \\
ECO               & none (inference-time prompt corruption) & prompt classifier \\
\textbf{Ours}     & \textbf{input embedding only}           & none \\
\bottomrule
\end{tabular}
\end{table}

\noindent
Our contributions are (i) the structural negative result of
Sec.~\ref{sec:negative}, (ii) subject-keyed input-embedding unlearning with
combinatorial locality, and (iii) taking the sensitive readout direction in
closed form from the LM-head row rather than estimating it. We explicitly do
\emph{not} claim the invariance regularizer as a contribution.

% ===========================================================================
\section{Limitations}
% ===========================================================================

\paragraph{$\Spe$ is bimodal.} Five of ten seeds sit at or above base and five
collapse to $6.3$--$9.5$, with none near the mean; the $\pm 2.32$ SD describes
a discrete cause rather than a spread. The mechanism is the same one that
breaks $\Spe$ at $m = 10$: with full coverage the edit includes \emph{common}
subject tokens, which neighborhood prompts do contain, so whether a seed's
draw includes such subjects determines its cluster. This is testable from
training-visible data alone by correlating each seed's $\Spe$ against the
corpus frequency of its edited rows. A per-row norm cap scaled by token
frequency is the natural remedy and is not yet evaluated.

\paragraph{Scope.} Results are reported for one benchmark and one model. The
method assumes the subject appears lexically in the prompts to be forgotten; it
therefore does not address multi-hop settings in which a fact is reached only
through an earlier inference step, since the subject string need not appear at
all.

\paragraph{$\Gen$ is not zero.} $\Gen = 1.70 \pm 2.00$, with four of ten seeds
at exactly zero. Five distinct hypotheses aimed at closing the remainder were
tested and rejected; only token coverage produced an improvement.


\appendix
% ===========================================================================
\section{Implementation details}
% ===========================================================================

This appendix specifies the architecture exactly as implemented, including the
details that turned out to matter empirically. Each subsection names the
behaviour and, where relevant, the failure it prevents.

\subsection{Parameterization and what is frozen}

Let $E \in \mathbb{R}^{V \times d}$ be the input-embedding matrix and
$W \in \mathbb{R}^{V \times d}$ the LM head, with $V = 128{,}256$, $d = 3072$
for Llama-3.2-3B-Instruct ($28$ decoder blocks).

\begin{itemize}
\item \textbf{Trainable:} a dense delta $\Delta \in \mathbb{R}^{|S| \times d}$
      over the selected subject rows $S$. This is the \emph{only} tensor passed
      to the optimizer.
\item \textbf{Frozen:} all $28$ blocks, $W$ (explicitly
      \texttt{requires\_grad\_(False)}), and all rows of $E$ outside $S$.
\item $\Delta$ is unconstrained in direction---no low-rank basis, no subspace
      projection. Restriction is unnecessary because locality does not come
      from restricting $\Delta$.
\end{itemize}

\paragraph{Untying.} $E$ and $W$ are untied before training and the identity
$E.\texttt{data\_ptr}() \neq W.\texttt{data\_ptr}()$ is asserted. With tied
weights, an edit to $E_t$ is simultaneously an edit to $W_t$, which changes
$t$'s logit in every context and voids the locality argument entirely.

\paragraph{Rows never edited.} $S$ is derived \emph{only} from the subject
string. The rows for \texttt{target\_true} and \texttt{target\_new} are never
modified in either $E$ or $W$; those tokens are used solely to index into the
logits when computing Eq.~\eqref{eq:loss}. This is the precise inversion of the
output-side family analysed in Sec.~\ref{sec:negative}.

\subsection{Subject row selection}

\paragraph{Two tokenizations.} The subject is tokenized both as
\texttt{" "\,+\,subject} (mid-sentence, as in the cloze prompt) and as
\texttt{subject} (sentence-initial, as in paraphrases). These differ under BPE,
and the union is taken. Special tokens are removed.

\paragraph{Liveness, split by role.} Candidate rows are intersected with the
tokens that actually occur in the record's own prompts, computed separately for
the direct prompt and for the synthetic paraphrases:
\begin{itemize}
\item \textbf{Direct-prompt liveness is mandatory.} Each record must retain at
      least one row occurring in its direct prompt, overriding the frequency
      threshold if every such token is common. $\Eff$ is scored there.
\item \textbf{Paraphrase-only rows} are optional extras, threshold-gated, kept
      for $\Gen$ coverage.
\item A record with no live direct-prompt token raises rather than training.
\end{itemize}
Pooling the two into one liveness test is insufficient: a token living only in
the sentence-initial form can pass a pooled test, train happily on synthetic
paraphrases, and never fire on the prompt $\Eff$ measures. In a real run this
left three records permanently unfixable while all rows appeared ``touched by
gradient''.

\paragraph{Frequency threshold.} A subject token that is also ordinary
vocabulary (\texttt{the}, \texttt{university}, \texttt{robert}) reintroduces
collateral damage. Rows may be filtered above a corpus-frequency threshold,
counted on the disjoint Wikipedia slice. In the reported configuration the
threshold is effectively disabled: filtering starved the edit
(Table~\ref{tab:ablation}).

\paragraph{Observed scale.} $236.3 \pm 11.5$ rows across $50$ records, i.e.\
$\approx 4.7$ rows per record, out of $128{,}256$.

\subsection{Applying the delta}

During training $\Delta$ is applied by a forward hook on the embedding module
rather than by mutating weights:
\begin{enumerate}
\item a lookup table maps vocabulary id $\rightarrow$ index into $S$, or $-1$;
\item for each input token, the row index is gathered and masked to zero where
      the token is not in $S$;
\item the masked correction is added to the module's output.
\end{enumerate}
This keeps $\Delta$ differentiable and confines the edit to positions where a
selected token actually appears. After training, $\Delta$ scaled by the
selected factor is written into $E$ with an in-place index-add, so the saved
artifact is an ordinary checkpoint requiring no hooks at inference.

\subsection{Training cases and the objective}

\paragraph{Case expansion.} Each record expands into one teacher-forced
next-token case per answer token. For answer tokens $(t_0,\dots,t_{k-1})$, case
$i$ has prompt $= \text{prompt} \Vert \mathrm{decode}(t_{<i})$ and target
$= \mathrm{decode}([t_i])$.

\paragraph{A bug worth recording.} An earlier port set the case target to the
\emph{whole} answer string at every position. Since the target-id helper
tokenizes that string and takes a single column, every position resolved to the
answer's \emph{first} token: for ``New York City'' all three positions
optimized ``New''. Single-token answers were coincidentally correct, so the
objective still appeared to converge while optimizing the wrong decisions. The
diagnostic that exposes it is an invariant now asserted in code: at scale
$\alpha = 0$ the materialized model \emph{is} the base model, so it must still
fail the objective; if scale $0$ reports zero failures, the training objective
is not scoring the decisions the evaluator scores.

\paragraph{Loss terms.} Eq.~\eqref{eq:loss} with $\lambda_m = 100$, $m = 6$,
$\lambda_2 = 10^{-4}$. Two further terms exist in the implementation and are
disabled in the reported configuration:
\begin{itemize}
\item a \emph{directional} term
      $\lVert \delta h - (\delta h^\top \hat u_s)\hat u_s \rVert^2$ confining
      the hidden-state change to
      $\hat u_s = W_{y^-}/\lVert W_{y^-}\rVert$, the closed-form readout
      direction for the sensitive answer (weight $0$; see
      Table~\ref{tab:ablation});
\item a \emph{context-invariance} term penalizing the across-context variance
      of the induced subject-representation shift (weight $0$; rejected).
\end{itemize}

\paragraph{On the $\ell_2$ term.} At $\lambda_2 = 10^{-4}$ and a typical delta
norm of $\approx 5$ over $223 \times 3072$ entries, this term contributes
$\approx 3.7\times10^{-9}$ against a margin term of order $10^{2}$. It is
therefore inert, and frequency-reweighting it has no measurable effect. Where a
per-row budget is wanted, the implementation instead offers a hard projection
after each optimizer step,
\begin{equation}
\Delta_i \leftarrow \Delta_i \cdot
\min\!\left(1, \frac{c_i}{\lVert \Delta_i \rVert}\right),
\qquad
c_i = \frac{c}{(1 + f_i)^{\alpha}},
\end{equation}
with $f_i$ the corpus frequency of row $i$. This binds regardless of loss
scaling. It is disabled ($c = 0$) in the reported result and is the proposed
remedy for the $\Spe$ bimodality.

\subsection{Optimization and scale selection}

Adam, learning rate $5\times10^{-4}$, weight decay $0$, batch size $4$, $1200$
steps, gradient-norm clipping at $1.0$. Non-finite losses or gradient norms
raise rather than being skipped.

The trained $\Delta$ is then evaluated at
$\alpha \in \{1, 0.875, \dots, 0.0078125, 0\}$ using the exact full-sequence
margins, and $\alpha^\star$ is chosen in three passes: minimize direct-only
failures; among those minimize combined direct$+$synthetic failures; among
remaining ties take the \emph{largest} $\alpha$. The final tie-break matters:
a naive smallest-on-tie rule collapses to $\alpha = 0$---a silent no-op---
whenever the combined objective cannot reach zero failures.

\subsection{Synthetic paraphrases}

Real paraphrases are held out, so training uses hand-authored substitutes:
three per record, drawn from a bank of two alternate cloze templates for each
of the $34$ MCF relation ids, optionally prefixed by an unrelated carrier
sentence sampled from the disjoint Wikipedia slice ($256$ sentences, $5$--$16$
words, deduplicated, deterministic under seed). Records whose relation id is
absent from the bank fall back to generic templates, and the count of such
records is reported so a missing relation id surfaces rather than silently
degrading.

Two facts about this component are worth stating plainly: synthetic failures
reached $0$ in every reported seed while real $\Gen$ did not, and in one
ablation synthetic failures \emph{improved} while real $\Gen$ \emph{worsened}.
The synthetic set is therefore not a reliable proxy for $\Gen$, and we do not
treat it as one.

\subsection{Recorded diagnostics}

Each run writes a configuration record containing, among others: the selected
token ids and per-record selection report (which rows were kept, dropped, or
absent from the prompts, with their frequencies); rows ever touched by
gradient; direct and synthetic failure counts and positions; the minimum
margin; the full scale sweep; the final delta norm; and the mean sensitive
readout $h^\top \hat u_s$ before and after editing, which provides mechanistic
evidence of forgetting independent of the NLL-based margin (observed drop
$7.71 \pm 0.48$).

\subsection{Variant for benchmarks without a replacement target}

MQuAKE and canonical ZsRE supply the sensitive answer only---their loaders
reject a neutral target---and score forgetting by argmax accuracy on that
answer rather than by pairwise preference. The architecture is unchanged; only
the reference in Eq.~\eqref{eq:loss} differs, becoming the highest-logit
non-sensitive token cached once on the \emph{base} model:
\begin{equation}
\mathcal{L} = \lambda_m\,
\mathrm{relu}\!\big(m - [\log p(y^{\text{comp}}\mid x) - \log p(y^{-}\mid x)]\big)
+ \lambda_2 \lVert\Delta\rVert^2 ,
\quad
y^{\text{comp}} = \arg\max_{t \neq y^{-}} \mathrm{logit}_t^{\text{base}} .
\end{equation}
This reference is stable by construction here: since only input-embedding rows
are edited, no output row changes and $y^{\text{comp}}$'s LM-head row is
identical to base throughout. Output-side editing methods must explicitly mask
their edited rows when selecting a competitor for exactly this reason.

\end{document}
